\documentclass[output=paper]{langscibook}
\ChapterDOI{10.5281/zenodo.22880296}
\title{\textit{Vavaatoa} – Kite fishing}
\author{Naphtali Maion}
\abstract{}
\begin{document}
\maketitle\label{chap:34}
\begin{figure}[H]
    \includegraphics[trim=0 50 0 0,clip,width=.9\linewidth]{figures/34_1.jpg}
    \caption{\textit{Vavaatoa} -- Kite fishing}
    \label{fig:34_1}
\end{figure}
\begin{xltabular}{\textwidth}{lXX}
	(\ref{ex:text27-1})&\textit{\textbf{Vagana vaa tea vavaatoa}}&\textbf{Fishing with the kite}\\
	(\ref{ex:text27-2})&\textit{A vavaatoa na vavaanaovo riori tea iravata.}&They let the kite fly in the wind.\\
	(\ref{ex:text27-3})&\textit{Evehee, saka ei haana komana vaan. Ahiki.}&But this is not in the village. No.\\
	(\ref{ex:text27-4})&\textit{Na dee riori tahii.}&They take (it) to the sea.\\
	(\ref{ex:text27-5})&\textit{A iravata to vavaanaovo vavaatoa suku riori saka hoa vabeera haana.}&The wind in which they let kites fly does not blow strongly.\\
	(\ref{ex:text27-6})&\textit{Bara o puhana to paku riori bona vagana bona saka beera haana mee.}&And the high tide during which they do this fishing is not strong either.\\
	(\ref{ex:text27-7})&\textit{A tahii toro ato koa bono takatakaku mooman.}&The sea water only needs to touch your ankles.\\
	(\ref{ex:text27-8})&\textit{A vavaatoa na vakis ni riori bona buaku a kanono ariono.}&They tie two \textit{ariono} strings to the kite.\\
	(\ref{ex:text27-9})&\textit{A peha ariono na vakis nana matana vavaatoa, a vaabuaku na vakis nana murinae.}&One \textit{ariono} string is tied to the front of the kite, the second one is tied behind it. \\
	(\ref{ex:text27-10})&\textit{A too vagana na ato kaakoo batana tea ariono vaa matana.}&The fisherman holds the string that is tied to its front.\\
	(\ref{ex:text27-11})&\textit{A matanima he na hana batana tea ariono vaa teo vamurina.}&The \textit{matanima}, the hook of the kite, hangs on the \textit{ariono} string in the back.\\
	(\ref{ex:text27-12})&\textit{Beori vaanaovo bono vavaatoa, saka vaanaovo pinopino haa riori.}&When they let the kites fly, they do not let them fly carelessly around.\\
	(\ref{ex:text27-13})&\textit{Na tara rori bona matanima be tei koa toonuna tahii, }&They see that the \textit{matanima} just stays on the surface of the sea, \\
	(\ref{ex:text27-14})&\textit{ore paa vaantee koa bari teebona.}&and then they let it stop there.\\
	(\ref{ex:text27-15})&\textit{A too vagana be pita mi bona vavaatoa, a matanima ore paa tereree bata koana toon na tahii.}&When the fisherman walks with the kite, the hook of the kite only hovers over the surface of the sea.\\
	(\ref{ex:text27-16})&\textit{O iana toro tara bari be tereree batana, ore paa kao bari,}&The fish will see it when it is hovering over the surface,\\
	(\ref{ex:text27-17})&\textit{ore paa kao bari ore paa no koto bari.}&then they approach and bite them.\\
	(\ref{ex:text27-18})&\textit{Pasi koto bata nom, beori hana.}&They will bite them, while they are hanging there.\\
	(\ref{ex:text27-19})&\textit{A vavaatoa na ponisi nana ore paa rahi gunaha bari, }&The kite is heavy, they pull it down, \\
	(\ref{ex:text27-20})&\textit{ore paa ato bona iana.}&and then they (the fishermen) grab the fish.\\
	(\ref{ex:text27-21})&\textit{Eve koa to paku bata vo riori, ore suguna komana }&That's how they do it until \\
	(\ref{ex:text27-22})&\textit{beo iana bero vaantee bono rake teori, ore paa voosu vahaa.}&they have as many fish as they wanted, and then they return home.\\
	(\ref{ex:text27-23})&\textit{Evei o iana to rahirahi maa riori tea vavaatoa,}&These fish that they catch with the kite,\\
	(\ref{ex:text27-24})&\textit{saka kavara rio iana haana, ahiki.}&it is not all fishes, no.\\
	(\ref{ex:text27-25})&\textit{O iana vai koa (1) hori, (2) hoohoa, (3) pasuaromo, (4) ponpon.}&It is only the following fishes: (1) \textit{hori} squirrel fish, (2) \textit{hoohoa} needle fish, (3) \textit{pasuaromo} emperor, (4) \textit{ponpon} needle fish.\\
	(\ref{ex:text27-26})&\textit{Bea tahii masi, ore paa vakavara vagana.}&When the tide goes down, they finish kite fishing.\\
\end{xltabular}
\begin{figure}[H]
    \captionsetup[subfigure]{oneside,margin={0cm,0cm}}
    \subfloat[]{\includegraphics[trim=0 0 0 0,clip,width=.45\linewidth]{figures/34_2.png}\label{fig:34_2}}\hfil
    \subfloat[]{\includegraphics[trim=0 0 0 0,clip,width=.45\linewidth]{figures/34_4.png}\label{fig:34_4}}
    \caption{(a) \textit{Hori} -- squirrel fish, (b) \textit{Pasuaromo} -- Emperor}
\end{figure}
\begin{figure}[H]
    \includegraphics[trim=0 0 0 0,clip,width=.7\linewidth]{figures/34_3.png}
    \caption{\textit{Hoohoa} -- needle fish}
    \label{fig:34_3}
\end{figure}
\begin{figure}[H]
    \includegraphics[trim=0 0 0 0,clip,width=.7\linewidth]{figures/34_5.png}
    \caption{\textit{Ponpon} -- needle fish}
    \label{fig:34_5}
\end{figure}
\end{document}